%% kapa: exam=1 round=27 year=2016 period=1 form=A subject=civil_law subject_ko=민법 seq=27 number=27
%% source: Q-Net 감정평가사 기출문제 (cst00309) artlSeq=5169402 file=extracted/1008/5169402/제27회 감정평가사 1차 문제/제27회 감정평가사 1교시문제(A형).pdf
%% conversion: kordoc 4.24.51 -> tools/gen1.py (auto); review=proofread-sonnet; flags=none
\begin{question}{27}
부동산 등기에 관한 설명으로 옳지 않은 것은? (다툼이 있으면 판례에 따름)
\begin{choices}
  \item 토지소유자가 그 지상에 건물을 신축하는 경우, 보존등기를 하여야 건물의 소유권을 취득한다.
  \item 무효인 중복등기에 바탕을 둔 등기부취득시효는 인정되지 않는다.
  \item 무효등기를 유용하는 합의는 그 합의 전에 등기상의 이해관계 있는 제3자가 없는 경우에는 유효하다.
  \item 증여에 의하여 부동산을 취득하였지만 등기원인을 매매로 기재하였더라도 그 등기는 유효하다.
  \item 「민법」에서는 등기의 추정력에 관한 규정을 두고 있지 않다.
\end{choices}
\end{question}
\answer{1}
